\section{Problem Formulation}

An inference scenario is a tuple
\begin{equation}
S=\langle M,W,H,R,C,\mathbf o,B_m\rangle,
\end{equation}
where $M$ identifies the model and allowable precisions; $W$ is a request
trace or distribution over arrivals and token lengths; $H$ describes the
target GPU node and power controls; $R$ is the serving runtime; $C$ contains
memory, administrative, model-quality, and service-level constraints;
$\mathbf o$ specifies the objectives; and $B_m$ is the real-measurement budget
in GPU-hours.

The candidate space is
\begin{equation}
\mathcal X=\mathcal B\!\times\!\mathcal K\!\times\!\mathcal P\!\times\!
\mathcal F,
\end{equation}
covering sequence and batched-token budgets ($\mathcal B$), KV-cache policy
($\mathcal K$), prefill and batching policy ($\mathcal P$), and fixed or
tunable precision/power controls ($\mathcal F$).  A deterministic compiler
produces the scenario-feasible pool $\mathcal X_S$ by removing candidates that
violate memory capacity, runtime support, or site policy.

For $m$ objectives and evidence stages
$e\in\mathcal E=\{\textsc{Sandbox},\textsc{Probe},\textsc{Verify}\}$, let
$\mathbf f^{(e)}(x)\in\mathbb R^m$ denote the noisy response obtained by
querying candidate $x$ at stage $e$.  \textsc{Verify} executes the full target
workload on the target GPU and is the reference fidelity.  Our primary
minimization vector is
\begin{equation}
\mathbf f^{(v)}(x)=
\langle E_{\rm tok},T_{\rm TTFT},T_{\rm TPOT},-\Theta\rangle,
\end{equation}
where latency entries denote scenario-specified tail statistics and $\Theta$
is goodput or throughput.  If precision or model variants are part of the
search, their quality metric is enforced as a constraint rather than traded
silently for energy.  Let $q_S(x)=1$ when the verified response distribution
satisfies every service and quality constraint in $C$, and let
$\mathbf u\prec\mathbf v$ denote Pareto dominance under minimization.  The
unknown target set is
\begin{align}
\mathcal P_S^*=\{x\in\mathcal X_S:
&q_S(x)=1,\ \nexists x'\in\mathcal X_S\ \text{s.t.}\nonumber\\[-2pt]
&q_S(x')=1\land
\mathbf f^{(v)}(x')\prec\mathbf f^{(v)}(x)\}.
\label{eq:true-frontier}
\end{align}

At step $t$, the controller has evidence
$\mathcal D_t=\{(x_i,e_i,y_i,c_i,z_i)\}_{i=1}^t$, where $y_i$ is a
measured or simulated outcome, $c_i$ is realized GPU-hour cost, and $z_i$
records execution status and provenance.  An action $a_t=(x_t,e_t)$ yields an
observation and updates a belief over every verified response
$\mathbf f^{(v)}(x)$.  The returned set $\widehat{\mathcal P}_T$ contains only
fully verified candidates that satisfy the measured SLO and are nondominated
among the verified candidates; it is an estimate of, not an assumed subset
of, $\mathcal P_S^*$.

For a search policy $\pi$, we minimize a terminal decision loss
\begin{equation}
\min_{\pi}\
\mathbb E\!\left[\mathcal L
(\widehat{\mathcal P}_T,\mathcal P_S^*)\right]
\quad\text{s.t.}\quad
\sum_{t=1}^{T}c(a_t)\leq B_m,
\label{eq:budgeted-objective}
\end{equation}
where $\mathcal L$ combines normalized hypervolume regret with penalties for
false SLO feasibility and unverified release.  In bounded replay spaces we
compute $\mathcal P_S^*$ exhaustively; on live spaces we report verified
hypervolume, feasible precision/recall over the measured pool, and independent
holdout verification.  All search curves use cumulative real GPU-hours because
a CPU query, short probe, and full target run are not equivalent observations.
